\section{Reconstituição digital das varreduras de MA-XRF}
\label{subsec:costura-hdf5}

Em algumas obras, a cobertura da região analisada foi obtida por meio de duas ou mais varreduras parciais com áreas de sobreposição. Essa divisão decorreu das dimensões das pinturas e das condições operacionais de aquisição. A reconstituição dos cubos foi realizada por uma rotina desenvolvida neste trabalho em Python, que compreendeu o alinhamento geométrico dos blocos, a equalização radiométrica relativa e a composição da região sobreposta por uma linha de costura. A rotina operou sobre dois blocos por vez, permitindo a reconstituição sucessiva dos cubos quando eram necessárias mais de duas varreduras.

Para o alinhamento espacial dos cubos armazenados em arquivos HDF5, foi gerada, para cada bloco, uma imagem de registro pela soma das contagens ao longo do eixo espectral. A rotina examinou faixas localizadas nas bordas adjacentes dos blocos e avaliou diferentes extensões candidatas de sobreposição. A escolha considerou a correlação normalizada entre essas regiões, com um termo adicional que favoreceu sobreposições mais extensas. Foram estimados o número de linhas ou colunas sobrepostas e o deslocamento espacial residual na direção transversal à costura. Os deslocamentos foram aplicados como translações inteiras dos blocos, sem interpolação espacial dos espectros.

A equalização radiométrica relativa empregou a região de sobreposição para estimar um fator multiplicativo entre os blocos. O uso de regiões sobrepostas para estimar ganhos e compensar diferenças radiométricas entre aquisições adjacentes é discutido por \citeonline{ban2022} no contexto da construção de mosaicos de imagens. Neste trabalho, esse princípio foi aplicado aos cubos espectrais de MA-XRF por meio do procedimento descrito a seguir.

Uma estimativa inicial do fator foi obtida a partir das faixas utilizadas no registro e posteriormente refinada na região de interseção dos cubos alinhados. Nesse refinamento, foi utilizada uma amostra regular de pares de pixels espacialmente correspondentes. Para cada pixel, as contagens foram somadas ao longo do eixo espectral, obtendo-se sua intensidade espectral total. Foram considerados apenas os pares nos quais ambas as intensidades eram finitas e superiores a zero. Em seguida, calculou-se a razão entre a intensidade do primeiro e a do segundo bloco. Quando havia pelo menos 50 razões válidas, foram desconsideradas aquelas inferiores ao 1º percentil e superiores ao 99º percentil. O fator de equalização foi definido pela mediana das razões consideradas. O bloco adotado como referência não recebeu reescalamento, enquanto o outro foi multiplicado pelo fator necessário para compatibilizar as escalas de intensidade. O mesmo fator foi aplicado a todos os canais espectrais de cada bloco. Na gravação em formato inteiro sem sinal de 32 bits (\texttt{uint32}), os valores reescalados foram arredondados.

Antes da composição do cubo final, a área útil foi delimitada pela interseção das extensões dos blocos na direção transversal à costura. Na costura vertical, o recorte foi realizado nas bordas laterais; na horizontal, nas bordas superior e inferior. Esse recorte excluiu as faixas periféricas fora da extensão comum, incluindo eventuais pixels válidos de apenas um dos blocos.

Na região de sobreposição, a origem dos espectros mantidos no cubo final foi definida por uma linha de costura, ou \textit{seam}, calculada por programação dinâmica. Seu traçado considerou as diferenças locais entre as intensidades espectrais integradas dos blocos e a continuidade espacial da linha de junção, sendo posteriormente suavizado por uma mediana móvel. Assim, diferentes partes da região sobreposta puderam ser preenchidas com dados de blocos distintos, mas cada posição do cubo final foi representada por um espectro proveniente de apenas uma aquisição, após a equalização, com o descarte do espectro redundante. Esse procedimento evitou a geração de espectros resultantes da combinação entre duas aquisições.

Os cubos reconstituídos foram utilizados na obtenção dos mapas elementares, na decomposição por NMF-KL e na posterior construção dos descritores empregados na PCA. As obras submetidas a esse procedimento e o número de varreduras parciais utilizadas na reconstituição dos respectivos cubos são apresentados na Tabela~\ref{tab:varreduras-costura}.
